% Updated B-SMART implementation chapter: multi-horizon RQ1 design, 7 November 2026.
\section{Implementation Scope and Research--Software Boundary}

B-SMART contains a multi-organisation operational web application and a separate offline empirical research workflow. The operational software supports organisation and branch configuration, product catalogue, sales, inventory movements, purchasing, finance, role-scoped analytics, and a recommendation review/decision log. The experiments in Chapter~\ref{chapter_experimentalResults} use a reproducible Python evaluation pipeline. The trained research models are \emph{offline artefacts}; they are not verified production inference endpoints. A page labelled Forecast or Segments is not, by itself, evidence that the thesis models are served in production.

\section{Operational Software Implementation}
\subsection{Application Architecture}

The repository uses React and Vite for its frontend, FastAPI and SQLAlchemy for the Python API, and configurable SQLite/PostgreSQL relational storage. Tenant and branch scope is retained throughout the business data model. The backend enforces capabilities through role-based permission checks; frontend visibility is not treated as sufficient isolation. Operational event records, import events, and recommendation decisions are tracked through auditable routes.

\subsection{Business Modules and Outputs}

The operational workflow includes authenticated organisation/branch configuration; product catalogue and point-of-sale orders; inventory movements; supplier and purchasing records; sales returns; expense, payment and ledger reporting; and dashboard indicators. The application computes traceable operational KPIs where source evidence is adequate. Cost coverage is distinguished from revenue rather than assuming that missing costs represent zero. Reported KPIs require explicit time-window definitions and eligible transaction statuses.

\subsection{Recommendation and Owner Decision Contract}

An actionable candidate should contain the decision cutoff, eligible source snapshot, suggested action, estimated benefit and uncertainty where supported, stock/budget/lead-time/capacity/consent constraints, and a reason for abstaining when essential data are absent. The owner can accept, reject, modify, or defer a recommendation; later outcome capture is a separate event. A predicted inactivity score is not a causal campaign-uplift estimate, and a demand forecast is not proof that stock was available to meet demand.

\section{Reproducible Offline Model Implementation}
\subsection{Data Validation and Empirical Scope}

The forecasting panel contains 561,152 complete shop--category--day rows: 32 shops, 16 categories, and 1,096 calendar days from 1 January 2023 to 31 December 2025. The panel has 512 shop--category series, no duplicate keys, and no missing shop--category--date combinations. The restricted reconstruction joined 736,683 sales lines to 365,447 orders and the product catalogue; the selected source export had no missing order or SKU references and no shop-ID mismatches. Three similarly named sales-line exports overlap and are not independent partitions; only one non-overlapping source copy was used for reconstruction.

All 32 shops in the supplied catalogue are labelled grocery/FMCG. The empirical results therefore concern this retail setting, not a representative sample of all SME sectors. Merchant origin, collection permissions, and the original extraction process have not been independently certified. Raw customer-level records and row-level predictions are excluded from this public repository.

The RQ1 wording retains its intended branch--SKU decision focus. The current reliable daily panel supports shop--category targets. SKU-level diagnostics show 13,271 observed shop--SKU pairs, an average of 52.7 selling days per pair over 1,096 days, and a 95.2\% daily zero rate. Accordingly, the executed primary panel experiment is shop--category; a fully validated branch--SKU forecast and SKU-allocation rule remain separate work.

\subsection{Daily Panel Features and Model Selection}

The daily model uses lagged target history and calendar context. The history-rich specification contains 75 predictors: target lags up to 364 days; trailing means, standard deviations and zero-demand rates over 7--84 days; positive-demand size summaries; recent-to-longer-term demand ratios; day-of-week, month and annual-cycle terms; and lagged shop, category and network demand context. All target-derived features are shifted so that the prediction for date \(t\) uses no sales from date \(t\) or later.

A global LightGBM model with an L1 objective is compared with the original smaller-feature LightGBM L2 model. L1 aligns the fitting loss with the absolute-error numerator of WAPE, but it can estimate a conditional median and understate expected volume. Therefore WAPE and aggregate forecast bias are reported together; the L1 daily model is not automatically selected as an inventory order-quantity model.

Transaction-derived lagged price, unit-cost, discount, gross-margin, phone-order-share and SKU-mix summaries were tested as an ablation. They did not improve daily WAPE in the first two rolling folds, so they are not included in the selected primary feature set. The finding does not imply that future-known promotions or market variables are unimportant; those variables were not available as reliable prospective inputs in the supplied panel.

\subsection{Direct Multi-Horizon Demand Forecasting}

To align forecasting with replenishment decisions, separate direct models predict cumulative observed demand over the next 7, 14 and 28 days. Each origin uses up to 84 days of past shop--category demand, demand intermittency and positive-size summaries, calendar terms, and lagged shop/category/network context. The target is the sum of the following \(h\) days, where \(h\in\{7,14,28\}\). Origins are spaced by \(h\) days, so evaluation target windows do not overlap within a horizon. This is a direct horizon-specific forecast, not a recursive sum of one-day predictions.

For each horizon, LightGBM L1 and Tweedie candidates are compared with a last-horizon repeat and a mean of the preceding four horizon blocks. Tweedie is retained as the primary inventory-planning candidate when its WAPE is close to L1 and its total-volume bias is lower. The L1 model remains a useful absolute-error benchmark. The horizon used in a later inventory policy must match the actual supplier lead time or review period; a 28-day forecast is not substituted for a 7-day lead-time forecast.

\subsection{Temporal and Held-Out-Shop Validation}

Three rolling-origin folds use target-end dates to prevent training labels from crossing a validation or test boundary. The test periods are 2025 Q1, Q2 and Q3; validation and calibration windows precede each test period. The October--December 2025 daily test used in earlier work has already been inspected and is not treated as a pristine final holdout.

A separate four-fold held-out-shop experiment excludes eight shops from model fitting in each fold and removes the shop identifier from the predictors. Historical sales from each held-out shop remain available as lagged context. This tests transfer to an unseen shop \emph{with historical sales available}; it is not zero-shot cold-start prediction and it does not establish transfer to a new sector.

Uncertainty was explored with 10th and 90th conditional-quantile models. Their empirical coverage was below the nominal 80\% interval target in the examined test folds. The intervals are therefore diagnostic and are not presented as calibrated service-level guarantees. Further temporal calibration is required before using an upper quantile as a stock-safety rule.

\section{RQ2, RQ3 and RQ4 Offline Implementation Boundaries}

RQ2 predicts no qualifying customer purchase during the following 90 days. Historical snapshots with overlapping outcome windows are excluded through temporal embargo. Stored Logistic Regression and LightGBM models are analysed offline. Aggregate TreeSHAP contributions are reported in raw-score (log-odds) space and do not prove causality or owner understanding.

RQ3 currently evaluates 2,500 synthetic-oracle scenarios. The new 7-, 14- and 28-day forecasts provide candidate demand inputs for a future policy replay, but they have not yet been connected to verified historical stock balances, receipts, supplier lead times, or realised purchase outcomes. The scenario experiment must not be described as a real-data backtest.

RQ4 currently audits 1,800 structured explanation cases. A supplied evidence-complete flag is not independent evidence verification. No SME-owner comprehension or usability study has been completed. Technical feature attribution, a structured business reason, and a human participant's understanding remain distinct outcomes.
